% SPDX-FileCopyrightText: Copyright (c) 2016-2026 Objectionary.com
% SPDX-License-Identifier: MIT

\begin{proof}[Proof of \cref{thm:dataization-confluent}]
Dataization is specified by the syntax-directed inference system of
  \cref{fig:dataization}, so confluence amounts to the relation $\Downarrow$
  being single-valued: whenever \(\phinoDataize{n}{e}{s_1}{\delta}{s_2}\)
  and \(\phinoDataize{n}{e}{s_1}{\delta'}{s_2'}\), we have \(\delta = \delta'\)
  and \(s_2 = s_2'\).
We argue by induction on the first derivation, after observing that at most
  one rule applies to any configuration \(\langle n, e, s_1 \rangle\),
  so both derivations are headed by the same rule.

The conclusions of the five rules partition every configuration, and $T$ falls
  under none of them.
If $n = T$, no rule matches, so dataization is undefined and there is nothing
  to disprove.
If \(n\) is neither a formation nor $T$, only \nameref{r:norm} matches.
If $n = [[ B ]]$ is a formation, the asset and attribute it carries select
  a single rule:
  \begin{inparaenum}[a)]
  \item a $D$-asset selects \nameref{r:delta} and an $L$-asset selects
    \nameref{r:fire}, the two being mutually exclusive in a normal form;
  \item with neither asset present, an $@$-attribute selects \nameref{r:box}
    and its absence selects \nameref{r:none},
  \end{inparaenum}
  the latter pair separated by the explicit conditions
  $[D, L] \cap B = \emptyset$ and $[D, L, @] \cap B = \emptyset$.
Hence no configuration matches two rules.

Because the same rule heads both derivations, it suffices to check that each
  rule maps its premises to a single result.
\nameref{r:delta} is an axiom whose conclusion reads the result directly off
  \(n\)---the attached data---leaving the state untouched, so it admits no other
  outcome.
\nameref{r:none} reduces to the dataization of $T$, which is undefined, so it
  yields no result and hence no second one.
\nameref{r:box} chains three premises: contextualization is a total function
  (\cref{sec:contextualization}) and yields a unique \(n\); normalization is
  confluent (\cref{sec:normalization}) and yields a unique normal form \(n_1\);
  and the dataization of \(n_1\) is unique by the induction hypothesis.
\nameref{r:norm} chains morphing, which is confluent (\cref{sec:morphing}) and
  yields a unique formation or $T$, with a dataization premise that is unique
  by the induction hypothesis.
\nameref{r:fire} evaluates the function \(F\) attached to the $L$-asset.
By \cref{sec:evaluation} a function is a total mapping that deterministically
  maps formation, universe, and state to an expression and a new state;
  being pure, its result depends on nothing else.
The universe \(e\) is fixed throughout the derivation and the state propagates
  only through the explicit threading of the rules, so both derivations present
  the call with identical arguments and it returns the same
  \(\langle n, s_2 \rangle\); normalization then yields a unique \(n_1\) and the
  dataization premise is unique by the induction hypothesis.
In every case the preceding premises are determined before the recursive one,
  so the sub-derivations dataize the same normal form in the same state and the
  induction hypothesis applies.
The two derivations therefore agree on the result, and dataization relates each
  normal form to at most one data and state.
\end{proof}
